\documentclass{article}
\usepackage{graphicx}
\usepackage{geometry}
\usepackage{pgfgantt}

\newcommand{\todo}[1]{\textcolor{red}{\textbf{TODO:} #1}}

\geometry{a4paper, margin=1in}
\title{Comparative Analysis of zk-SNARK
Instantiations for Ensuring Ballot Validity: Structure and Key Points}
\author{}
\date{}

\begin{document}

\maketitle

\section*{Abstract}

\section{Introduction}
\begin{itemize}
    \item Introduce E-voting
    \item Introduce concept of proving ballot validity and propose GPZKPs as a possible solution.
    \item Motivate Implementation in Circom
\end{itemize}

\section{Related Work}
\begin{itemize}
    \item Reference your paper (zk-SNARKs for ballot validity: A feasibility study) as starting point for this thesis
    \item Reference Kryvos, MoonMathManual, ... (see Zotero)
\end{itemize}

!!!POTENTIALLY STARTING HERE INTO METHODOLOGY!!!
\section{Voting Systems}
\begin{itemize}
    \item Definition Voter, Voting Authority, Voting System
    \item Outline problems with this approach (Problem: Proof for ballot validity is missing)
\end{itemize}

\section{Proof Systems}
\subsection{Definition}
\subsection{ZKPs, GPZKPs, SNARKs}
\subsection{Algebraic circuits, R1Cs and QAP}
\subsection{Groth16}
\subsection{Circom}
\begin{itemize}
    \item Just overall purpose and usage, not any of the implementations here
    \item Maybe small example circuit
\end{itemize}

\section{snarkjs}
\begin{itemize}
    \item Usage together with Circom
\end{itemize}

\section{Zero Knowledge Voting Systems}
\todo{Maybe move Voting Systems here?}
\begin{itemize}
    \item Show how ZKPs can be used to prove ballot validity
    \item Definition ZK Voter, ZK Voting Authority and ZK Voting System
\end{itemize}
\subsection{Proving Ballot Validity}
\begin{itemize}
    \item Argue that we need to show two statements:
    \begin{enumerate}
        \item The plain ballot is from the allowed choice space.
        \item The encrypted ballot belongs to the chosen plain ballot.
    \end{enumerate}
    \item Introduce high-level circuits for the two statements
    \item Combine the two individual circuits into one.
\end{itemize}

\section{Exponential ElGamal Encryption}
\begin{itemize}
    \item Definition
    \item Motivate why choosing elliptic curves as the groups for EEG makes sense
    \item Maybe short comparison with PVC (or reference zk-SNARKs for ... for that)
\end{itemize}

\section{Elliptic Curves in Circom}
\begin{itemize}
    \item Motivate why we use Montgomery curves
\end{itemize}
\subsection{The Group Law}
\subsection{Scalar Multiplication}
\subsubsection{Montgomery Ladder}
\subsubsection{y-Recovery}

\section{Our Voting systems}\todo{Better title}\\
For each of election type:
\begin{itemize}
    \item Choice space definition
    \item Circuit for Asserting Ballot Validity in Circom (Potentially for small example if to complicated otherwise)
\end{itemize}
\subsection{Single-Vote}
\subsection{Multi-Vote}
\subsubsection{Multi-Vote with Rules (MWR)}
\subsection{Line-Vote}
\subsection{Borda Election types}
\subsubsection{Pointlist-Borda}
\subsubsection{Borda Tournament Style (BTS)}
\subsection{Condorcet Methods}

\section{Benchmarks}
\begin{itemize}
    \item Test Suite
    \item Benchmarked System
\end{itemize}
!!!POTENTIALLY ENDING HERE INTO METHODOLOGY!!!

\section{Results}

\section{Discussion}

\section{Conclusion}

\section{Future Work}
\begin{itemize}
    \item Other Groth-16 SNARK backends such as gnark
    \item Other GPZKPs usch as PLONK, FFLONK, \dots
    \item (Twisted Edwards Curves)
    \item Other backend Curves (e.g., BLS12-381) (I might still try that If I have time)
\end{itemize}

\section{Acknowledgements}

\section*{References}

\section*{Appendix}

\end{document}